\documentclass{article}
\usepackage[paperwidth=60cm,paperheight=120cm,margin=0pt]{geometry}
\usepackage{fontspec}\setmainfont{Latin Modern Roman}
\usepackage{xeCJK}\setCJKmainfont{Microsoft YaHei}
\usepackage{amsmath, amssymb, siunitx}
\usepackage[american]{circuitikz}
\ctikzset{bipoles/length=1.2cm}
% 不用 standalone：它自家的裁边在 xelatex 下不生效，长图会被 letter 页面截掉尾部。
% 改用 preview 把页面尺寸严格设成内容尺寸（自动裁边）。
\usepackage[active,tightpage]{preview}
\PreviewEnvironment{circuitikz}
\setlength\PreviewBorder{8pt}
\pagestyle{empty}
\begin{document}

\begin{circuitikz}
\small
\node[anchor=east, font=\bfseries\small] at (4.55,0.00) {MCU\ /\ 电源};
\node[anchor=west, font=\bfseries\small] at (4.95,0.00) {网络名};
\node[anchor=west, font=\bfseries\small] at (12.45,0.00) {连接器针脚};
\node[anchor=west, font=\bfseries\small] at (12.30,-0.60) {HC-06 蓝牙（XH2.54-4P）};
\node[anchor=west, font=\scriptsize\itshape] at (12.30,-0.98) {仅调试用；比赛时不接（全自动、禁遥控）};
\node[anchor=east, font=\small] at (4.55,-1.84) {5V\ \textperiodcentered\ 裸模块为 3.3V};
\draw (4.85,-1.84) -- (12.30,-1.84);
\node[anchor=south, font=\scriptsize\ttfamily] at (8.57,-1.79) {5V};
\node[anchor=west, font=\small] at (12.48,-1.84) {1\,--\,VCC};
\node[anchor=east, font=\small] at (4.55,-2.58) {GND};
\draw (4.85,-2.58) -- (12.30,-2.58);
\node[anchor=south, font=\scriptsize\ttfamily] at (8.57,-2.53) {GND};
\node[anchor=west, font=\small] at (12.48,-2.58) {2\,--\,GND};
\node[anchor=east, font=\small] at (4.55,-3.32) {PA3\ \textperiodcentered\ USART2\_RX};
\draw (4.85,-3.32) -- (12.30,-3.32);
\node[anchor=south, font=\scriptsize\ttfamily] at (8.57,-3.27) {BT\_TX};
\node[anchor=west, font=\small] at (12.48,-3.32) {3\,--\,TX};
\node[anchor=east, font=\small] at (4.55,-4.06) {PA2\ \textperiodcentered\ USART2\_TX};
\draw (4.85,-4.06) -- (12.30,-4.06);
\node[anchor=south, font=\scriptsize\ttfamily] at (8.57,-4.01) {BT\_RX};
\node[anchor=west, font=\small] at (12.48,-4.06) {4\,--\,RX};
\draw[rounded corners=2pt] (12.30,-1.42) rectangle (18.40,-4.36);
\node[anchor=west, font=\bfseries\small] at (12.30,-4.96) {遥测串口 USB-TTL（XH2.54-4P）};
\node[anchor=west, font=\scriptsize\itshape] at (12.30,-5.34) {接笔记本，115200 8N1（后续可升 460800）};
\node[anchor=east, font=\small] at (4.55,-6.20) {5V\ \textperiodcentered\ TTL 电源脚可不接};
\draw (4.85,-6.20) -- (12.30,-6.20);
\node[anchor=south, font=\scriptsize\ttfamily] at (8.57,-6.15) {5V};
\node[anchor=west, font=\small] at (12.48,-6.20) {1\,--\,5V};
\node[anchor=east, font=\small] at (4.55,-6.94) {PA9\ \textperiodcentered\ USART1\_TX};
\draw (4.85,-6.94) -- (12.30,-6.94);
\node[anchor=south, font=\scriptsize\ttfamily] at (8.57,-6.89) {UART1\_TX};
\node[anchor=west, font=\small] at (12.48,-6.94) {2\,--\,RXD};
\node[anchor=east, font=\small] at (4.55,-7.68) {PA10\ \textperiodcentered\ USART1\_RX};
\draw (4.85,-7.68) -- (12.30,-7.68);
\node[anchor=south, font=\scriptsize\ttfamily] at (8.57,-7.63) {UART1\_RX};
\node[anchor=west, font=\small] at (12.48,-7.68) {3\,--\,TXD};
\node[anchor=east, font=\small] at (4.55,-8.42) {GND};
\draw (4.85,-8.42) -- (12.30,-8.42);
\node[anchor=south, font=\scriptsize\ttfamily] at (8.57,-8.37) {GND};
\node[anchor=west, font=\small] at (12.48,-8.42) {4\,--\,GND};
\draw[rounded corners=2pt] (12.30,-5.78) rectangle (18.40,-8.72);
\node[anchor=west, font=\bfseries\small] at (12.30,-9.32) {SWD 调试排针（1×5）};
\node[anchor=west, font=\scriptsize\itshape] at (12.30,-9.70) {NRST 必须引出：关掉 JTAG 后靠它连接下复位救板};
\node[anchor=east, font=\small] at (4.55,-10.56) {3V3\ \textperiodcentered\ 来自核心板};
\draw (4.85,-10.56) -- (12.30,-10.56);
\node[anchor=south, font=\scriptsize\ttfamily] at (8.57,-10.51) {3V3};
\node[anchor=west, font=\small] at (12.48,-10.56) {1\,--\,3V3};
\node[anchor=east, font=\small] at (4.55,-11.30) {PA13};
\draw (4.85,-11.30) -- (12.30,-11.30);
\node[anchor=south, font=\scriptsize\ttfamily] at (8.57,-11.25) {SWDIO};
\node[anchor=west, font=\small] at (12.48,-11.30) {2\,--\,SWDIO};
\node[anchor=east, font=\small] at (4.55,-12.04) {PA14};
\draw (4.85,-12.04) -- (12.30,-12.04);
\node[anchor=south, font=\scriptsize\ttfamily] at (8.57,-11.99) {SWCLK};
\node[anchor=west, font=\small] at (12.48,-12.04) {3\,--\,SWCLK};
\node[anchor=east, font=\small] at (4.55,-12.78) {NRST\ \textperiodcentered\ 核心板 R 脚};
\draw (4.85,-12.78) -- (12.30,-12.78);
\node[anchor=south, font=\scriptsize\ttfamily] at (8.57,-12.73) {NRST};
\node[anchor=west, font=\small] at (12.48,-12.78) {4\,--\,NRST};
\node[anchor=east, font=\small] at (4.55,-13.52) {GND};
\draw (4.85,-13.52) -- (12.30,-13.52);
\node[anchor=south, font=\scriptsize\ttfamily] at (8.57,-13.47) {GND};
\node[anchor=west, font=\small] at (12.48,-13.52) {5\,--\,GND};
\draw[rounded corners=2pt] (12.30,-10.14) rectangle (18.40,-13.82);
\node[anchor=west, font=\bfseries\small] at (12.30,-14.42) {按键 / 蜂鸣器 / 指示灯};
\node[anchor=west, font=\scriptsize\itshape] at (12.30,-14.80) {核心板已自带复位键、BOOT0 跳线帽、电源灯 → PCB 上不再重复};
\node[anchor=east, font=\small] at (4.55,-15.66) {PB2\ \textperiodcentered\ 10k 上拉 + 100nF};
\draw (4.85,-15.66) -- (12.30,-15.66);
\node[anchor=south, font=\scriptsize\ttfamily] at (8.57,-15.61) {KEY};
\node[anchor=west, font=\small] at (12.48,-15.66) {—\,--\,用户按键 → GND};
\node[anchor=east, font=\small] at (4.55,-16.40) {PB1\ \textperiodcentered\ 定时器中断翻转方波};
\draw (4.85,-16.40) -- (12.30,-16.40);
\node[anchor=south, font=\scriptsize\ttfamily] at (8.57,-16.35) {BUZZER};
\node[anchor=west, font=\small] at (12.48,-16.40) {—\,--\,有源蜂鸣器 + NPN 三极管 + 1k 基极电阻};
\node[anchor=east, font=\small] at (4.55,-17.14) {12V\ \textperiodcentered\ 1k 限流};
\draw (4.85,-17.14) -- (12.30,-17.14);
\node[anchor=south, font=\scriptsize\ttfamily] at (8.57,-17.09) {LED\_12V};
\node[anchor=west, font=\small] at (12.48,-17.14) {—\,--\,12V 电源指示灯 → GND};
\draw[rounded corners=2pt] (12.30,-15.24) rectangle (18.40,-17.44);
\node[anchor=west, font=\bfseries\small] at (12.30,-18.04) {预留排针区（2×4）};
\node[anchor=west, font=\scriptsize\itshape] at (12.30,-18.42) {空闲 IO + 电源；示波器 / 逻辑分析仪 / 后续扩展。PA12 只作输出};
\node[anchor=east, font=\small] at (4.55,-19.28) {PB3\ \textperiodcentered\ 需关 JTAG};
\draw (4.85,-19.28) -- (12.30,-19.28);
\node[anchor=south, font=\scriptsize\ttfamily] at (8.57,-19.23) {SPARE1};
\node[anchor=west, font=\small] at (12.48,-19.28) {1\,--\,SPARE1};
\node[anchor=east, font=\small] at (4.55,-20.02) {PB4\ \textperiodcentered\ 需关 JTAG};
\draw (4.85,-20.02) -- (12.30,-20.02);
\node[anchor=south, font=\scriptsize\ttfamily] at (8.57,-19.97) {SPARE2};
\node[anchor=west, font=\small] at (12.48,-20.02) {2\,--\,SPARE2};
\node[anchor=east, font=\small] at (4.55,-20.76) {PB5\ \textperiodcentered\ 空闲};
\draw (4.85,-20.76) -- (12.30,-20.76);
\node[anchor=south, font=\scriptsize\ttfamily] at (8.57,-20.71) {SPARE3};
\node[anchor=west, font=\small] at (12.48,-20.76) {3\,--\,SPARE3};
\node[anchor=east, font=\small] at (4.55,-21.50) {PA12\ \textperiodcentered\ 原 USB DP，1.5k 上拉，勿作输入};
\draw (4.85,-21.50) -- (12.30,-21.50);
\node[anchor=south, font=\scriptsize\ttfamily] at (8.57,-21.45) {SPARE4};
\node[anchor=west, font=\small] at (12.48,-21.50) {4\,--\,SPARE4};
\node[anchor=east, font=\small] at (4.55,-22.24) {3V3\ \textperiodcentered\ 来自核心板};
\draw (4.85,-22.24) -- (12.30,-22.24);
\node[anchor=south, font=\scriptsize\ttfamily] at (8.57,-22.19) {3V3};
\node[anchor=west, font=\small] at (12.48,-22.24) {5\,--\,3V3};
\node[anchor=east, font=\small] at (4.55,-22.98) {5V};
\draw (4.85,-22.98) -- (12.30,-22.98);
\node[anchor=south, font=\scriptsize\ttfamily] at (8.57,-22.93) {5V};
\node[anchor=west, font=\small] at (12.48,-22.98) {6\,--\,5V};
\node[anchor=east, font=\small] at (4.55,-23.72) {GND};
\draw (4.85,-23.72) -- (12.30,-23.72);
\node[anchor=south, font=\scriptsize\ttfamily] at (8.57,-23.67) {GND};
\node[anchor=west, font=\small] at (12.48,-23.72) {7\,--\,GND};
\node[anchor=east, font=\small] at (4.55,-24.46) {12V\ \textperiodcentered\ 经 3A 保险丝};
\draw (4.85,-24.46) -- (12.30,-24.46);
\node[anchor=south, font=\scriptsize\ttfamily] at (8.57,-24.41) {12V};
\node[anchor=west, font=\small] at (12.48,-24.46) {8\,--\,12V};
\draw[rounded corners=2pt] (12.30,-18.86) rectangle (18.40,-24.76);
\node[anchor=west, font=\scriptsize\itshape] at (0,-25.66) {网络名 = 嘉立创 EDA 里的网络标号：两边同名即同一根线。电源（5V / 3V3 / GND / 12V）同样用网络标号表示，不画共用母线。逐针定义、电气约束与改动规则见 电气连接规范.md；引脚契约见 docs/引脚分配表.md。};
\end{circuitikz}
\end{document}
